\section*{Lernziele Kapitel 1: Erste Programme und Schleifen}
\begin{todolist}
\item Ich kenne die Befehle, die auf Seite 7 im Buch aufgelistet sind, und kann diese auch einsetzen.
\item Ich kann mit der Turtle einfache Formen zeichnen (Dreiecke, Rechtecke, Blitz, …)
\item Ich kann mit einem breiten Stift ausgefüllte Rechtecke zeichnen (Türe beim Haus)
\item Ich setze die \lstinline[style=TigerJython]|repeat|-Schleife ein, um sich wiederholende Tätigkeiten auszuführen.
\item Ich kann mit der Turtle regelmässige Vielecke zeichnen.
\item Ich kann ``Kreise'' als Vielecke mit grosser Eckenzahl zeichnen.
\item Ich kann \lstinline[style=TigerJython]|repeat|-Schleifen innerhalb von \lstinline[style=TigerJython]|repeat|-Schleifen einsetzen.
\item Ich kann Blütenblätter zeichnen, die aus Drittel, Viertel-, Fünftel-, usw.-Kreise bestehen.
\item Ich kann Text in der Konsole ausgeben.
\item Ich setze die Operationen \lstinline[style=TigerJython]|+|, \lstinline[style=TigerJython]|*| und \lstinline[style=TigerJython]|\n| ein, um komplizierte Strings (Zeichenketten) zu erzeugen.
\item Ich kenne die Operationen auf Zahlen, die auf Seite 19 oben aufgelistet sind, und kann diese auch einsetzen.
\item Ich kann innerhalb eines Programms eine Streckenlänge mit Pythagoras berechnen.
\end{todolist}
\section*{Lernziele Kapitel 2: Modularer Programmentwurf}
\begin{todolist}
\item Ich kenne die Befehle, die auf Seite 35 im Buch aufgelistet sind.
\item Ich kann eigene Befehle in Python definieren und einsetzen.
\item Ich achte darauf, dass am Ende eines Befehls, die Turtle den gleichen Zustand, wie am Anfang hat.
\item Ich nutze den modularen Programmentwurf, um komplexe Probleme in Teilprobleme aufzuteilen.
\item Ich kombiniere einfache Befehle, um komplexe Probleme zu lösen (Beispiel Häuserreihe).
\item Ich kann neue Befehle definieren, in denen andere eigene Befehle vorkommen (Beispiel 2.4).
\item Ich kann Animationen programmieren.
\end{todolist}

\section*{Weitere Lernziele (bis zu Kapitel 2)}
\begin{todolist}
\item Ich verstehe, wie sich die Turtle auf einem \texttt{x,y}-Koordinatennetz bewegt und kann sie mit den Befehlen \lstinline[style=TigerJython]|setPos(x,y)| und \lstinline[style=TigerJython]|heading(...)| positionieren und orientieren
\item Ich kann Formen farbig ausfüllen indem ich die Befehle \lstinline[style=TigerJython]|clear("...")|, \lstinline[style=TigerJython]|setFillColor("...")|, \lstinline[style=TigerJython]|startPath()| und \lstinline[style=TigerJython]|fillPath()| verwende.
\item Ich kann man Zufallszahlen mit der Library \lstinline[style=TigerJython]|random| generieren und verwenden, beispielsweise um die Position der Turtle zu ändern
\item Ich kann neue Farben generieren und verwenden (beispielsweise mit \lstinline[style=TigerJython]|setPenColor("...")| oder \lstinline[style=TigerJython]|setFillColor("...")|), indem ich RGB-Kodierungen wie \texttt{244, 12, 20} verwende.
\end{todolist}

\section*{Lernziele Kapitel 5: Listen}

\begin{todolist}
\item Ich kann Listen in Python erstellen.
\item Ich kann mit dem Index einzelne Elemente der Liste abrufen und verändern.
\item Ich kann die Anzahl Elemente einer Liste in Python berechnen lassen.
\item Ich kann die Elemente einer Liste in einer Schleife (\lstinline[style=TigerJython]|repeat|, \lstinline[style=TigerJython]|while| oder \lstinline[style=TigerJython]|for...in|) durchlaufen.
\item Ich kann ein Programm schreiben, das in einer Liste die Elemente mit bestimmten Eigenschaften findet (z.B. das Maximum oder alle ungeraden Zahlen).
\item Ich kann \textit{bubble sort} in Python programmieren.
\item Ich kann die binäre Suche in Python programmieren.
\item Ich verstehe, wie der Computer mit Listen arbeitet und kenne die Probleme, die bei \textit{copy by reference} entstehen.
\item Ich kann eine unabhängige Kopie einer bestehenden Liste erzeugen.
\item Ich kann Listen mit dem Befehl \lstinline[style=TigerJython]|.append(...)| erweitern.
\item Ich kann Listen mit dem Befehl \lstinline[style=TigerJython]|.pop()| kürzen.
\item Ich kann eine \lstinline[style=TigerJython]|for...in|-Schleife durch eine \lstinline[style=TigerJython]|while|-Schleife ersetzen.
\item Ich kann die Elemente einer Liste auf einen Wert reduzieren (z.B. auf die Summe der Listenelemente berechnen)
\end{todolist}

\section*{Lernziele Kapitel 7: Funktionen}
\begin{todolist}
\item Ich kann Funktionen in Python definieren, die einen oder mehrere Parameter haben.
\item Ich kann Funktionen aufrufen und in einem Ausdruck verwenden (z.B. \lstinline[style=TigerJython]|if rechteck_flaeche(19, 12) < 100: ...|)
\end{todolist}


